\begin{abstract}

    \begin{center}
    \large{Применение методов формальной верификации для анализа правил фильтрации сетевого трафика} \\[0.5 cm]
    \large{\textit{Дурнов Алексей Николаевич}} \\[1 cm]
    \end{center}

    Работа посвящена разработке системы автоматического анализа и верификации правил фильтрации сетевого трафика для выявления ошибок конфигурации сетевых устройств.
    В ходе работы был проведён обзор существующих методов анализа правил фильтрации; установлено, что известные подходы рассчитаны под конкретную платформу или ограниченный набор полей и не покрывают всё многообразие конфигураций современных межсетевых экранов.
    Для устранения этого ограничения была разработана унифицированная вендор-независимая модель представления правил фильтрации и объектов политики безопасности, охватывающая как классические списки контроля доступа, так и политики межсетевых экранов нового поколения, включая отечественные решения.
    За основу анализа взята теоретико-множественная модель FIREMAN, расширенная в рамках работы.
    Помимо классических аномалий между правилами фильтрации, анализатор обнаруживает аномалии внутри одного правила (выключенные, истёкшие и неиспользуемые правила) и ошибки в именованных объектах политики безопасности.
    Для правил с неизвестными анализатору условиями применяется приближённый анализ, исключающий ложные срабатывания.
    Разработанный анализатор был реализован и интегрирован в систему инвентаризации и валидации сетевых инфраструктур, разрабатываемую в ООО <<Дипси Нетворкс>>.
    Проведённое нагрузочное тестирование показало применимость анализатора к анализу больших наборов правил межсетевых экранов нового поколения.

\end{abstract}
\newpage
